\documentclass[11pt]{article}
\usepackage{fontspec}
\usepackage[french,english]{babel}
\usepackage{amsmath}
\usepackage{amssymb}
\usepackage{amsthm}
\usepackage{mathrsfs}
\usepackage{tikz-cd}
\usepackage{enumitem}

\newcommand{\printedpage}[1]{\marginpar{\footnotesize p.~#1}}
\newcommand{\msp}[1]{\marginpar{\footnotesize [#1]}}
\newcommand{\Filt}{\mathrm{Filt}}
\newcommand{\Fl}{\mathrm{Fl}}
\newcommand{\Ob}{\mathrm{Ob}}
\newcommand{\card}{\operatorname{card}}
\newcommand{\uHom}{\underline{\mathrm{Hom}}}
\newcommand{\uFl}{\underline{\mathrm{Fl}}}
\newcommand{\into}{\lhook\joinrel\longrightarrow}

\setlength{\parindent}{0pt}

\begin{document}

\printedpage{39}
% head: A. Grothendieck : Dérivateurs, Ch. XV, Édition des Universités de Montpellier II et Paris VII
% foot: Version 5 novembre 2010

\msp{page 67}
The equivalence of $(a)$ and $(a')$, or of $(b)$ and $(b')$ is more or less tautological. The equivalence of $(a)$ and $(b)$ follows at once from the N.B. included in $1^{\circ}$, taking into account that the cardinals $\pi$ such that the objects of $\Filt^{\pi}$ are $\pi$-accessible are cofinal in the set of cardinals in $\mathfrak{U}$. (\textit{Cf.} 3.2, page 16).

\textbf{Proposition 5.13.} \textit{Every finite union of L-$\pi$-accessible subsets of} $\Ob\mathcal{M}$ \textit{is L-$\pi$-accessible. The intersection of a family $(M_{\alpha})_{\alpha\in A}$ of L-$\pi$-accessible subsets, with} $\card A \leq \pi$\textit{, is L-$\pi$-accessible.}

\textsc{Proof}. The first assertion comes from the fact that a \textit{finite} union of L-stable subsets of $M$ is L-stable, and that we have likewise for the notion of $\pi$-accessibility. The two things are moreover proved practically in the same way, using the

\textbf{Lemma 5.14.} \textit{Let $I$ be a filtering ordered set, $(I_{\alpha})_{\alpha\in A}$ a} finite \textit{nonempty family of subsets of $I$, such that $\bigcup I_{\alpha} = I$. Then there exists $\alpha \in A$ such that $I_{\alpha}$ is cofinal.}

\msp{page 68}
\textsc{Proof} by contradiction. If for no $\alpha$, $I_{\alpha}$ were cofinal, we would conclude from this that for every $\alpha$ there exists an $i_{\alpha}$ such that $I/i_{\alpha} \cap I_{\alpha} = \emptyset$. Let $i$ be an upper bound for the $i_{\alpha}$, we will thus have $I/i \cap I_{\alpha} = \emptyset \ \ \forall\, \alpha \in A$, contrary to the hypothesis that $\bigcup I_{\alpha} = I$.

The case of an \textit{intersection} of L-$\pi$-accessible subsets: it is of course L-stable, so it remains to prove that it is $\pi$-accessible, \textit{i.e.} satisfies condition $\mathrm{b}_{\pi}$ of \textit{loc. cit.}

\msp{page 69}
Let then
\[
X \;=\; \varinjlim_{i\in I} X_{i} \; , \qquad\quad X_{i} \in \Filt^{\pi} \; ,
\]
be a $\pi$-adapted representation of $X$, with $I$ a $\pi$-adapted ordered set. Thus we have
\[
\begin{aligned}
X \in M_{\alpha} \quad &\Longleftrightarrow \quad \text{the set } I_{M_{\alpha}} \text{ of } i \in I \text{ such that } X_{i} \in M_{\alpha}\\
&\phantom{\Longleftrightarrow} \quad \text{is cofinal in } I.
\end{aligned}
\tag*{$(*)$}
\]
Moreover, since $M_{\alpha}^{\pi} = M_{\alpha} \cap \Filt^{\pi}$ is stable under filtering $\varinjlim$ of cardinality $\leq \pi$, we see at once that the subset $I_{M_{\alpha}}$ of $I$ is stable under suprema in $I$ of filtering subsets $J$ of $I_{M_{\alpha}}$ of cardinality $\leq \pi$. Thus $I_{M_{\alpha}}$ cofinal is equivalent to $I_{M_{\alpha}}$ $\pi$-\textit{cofinal} (\textit{cf.} 5.8, page 59). Let $M = \bigcap_{\alpha\in A} M_{\alpha}$, I want to prove that
\[
X \in M \quad \Longleftrightarrow \quad \text{the set } I_{M} \ (= \textstyle\bigcap I_{M_{\alpha}}) \text{ is cofinal in } I.
\tag*{$(**)$}
\]
Now $(*)$ implies that $X \in M$ $(= \bigcap M_{\alpha})$ is equivalent to $I_{M_{\alpha}}$ $\pi$-cofinal for every $\alpha$. So all in all we want to prove
\[
I_{M_{\alpha}} \text{ cofinal in } I \text{ for every } \alpha \quad\ \Longleftrightarrow \quad\ \textstyle\bigcap I_{M_{\alpha}} \text{ is cofinal in } I.
\tag*{$(***)$}
\]

\msp{page 70}
We of course have $\Longleftarrow$, so it remains to prove $\Longrightarrow$. If the index set $A$ is finite, this follows from Lemma 5.8, taking into account that the cofinality of the $I_{M_{\alpha}}$ implies their $\pi$-cofinality. We must therefore simply improve 5.8, by the

\printedpage{40}
% head: A. Grothendieck : Dérivateurs, Ch. XV, Édition des Universités de Montpellier II et Paris VII
% foot: Version 5 novembre 2010

\textbf{Lemma 5.15.} \textit{Let $\pi$ be a cardinal, $I$ a $\pi$-adapted filtering ordered set, and let $(I_{\alpha})_{\alpha\in A}$ be a family of subsets $I_{\alpha}$ that are $\pi$-cofinal} (cf. 5.8), \textit{such that} $\card A \leq \pi$. \textit{Then $\bigcap I_{\alpha}$ is $\pi$-cofinal.}

It suffices to prove that it is cofinal, for the property of stability under suprema of filtering subsets of cardinality $\leq \pi$ is obvious. Replacing $I$ by $I/i_{0}$, and the $I_{\alpha}$ by their traces on $I/i_{0}$ if need be, the question is simply to prove that $\bigcap_{\alpha} I_{\alpha} \neq \emptyset$. We are going to construct an element $s$ of $\bigcap I_{\alpha}$ by transfinite induction. Let us take a well-order on $A$, and prove by transfinite induction
\msp{page 71}
that for every $\alpha \in A$, the set
\[
J_{\alpha} \;=\; \bigcap_{\beta<\alpha} I_{\beta}
\]
is cofinal. If we know that $J_{\alpha}$ is cofinal, the same holds for $J_{\alpha+1} = J_{\alpha} \cap I_{\alpha}$, by 5.8. So we are reduced to proving that $J_{\alpha}$ is cofinal if $\alpha$ is a limit ordinal, thus to proving Lemma 5.15 for a \textit{decreasing} transfinite sequence of $\pi$-cofinal subsets $J_{\beta}$ $(\beta < \alpha)$ when $\card \alpha \leq \pi$, such that $J_{\beta} = \bigcap_{\beta'<\beta} J_{\beta'}$ when $\beta$ is a limit. To prove that the intersection is $\pi$-cofinal -- and we are reduced to proving that it is nonempty. We are going to construct by transfinite induction an \textit{increasing} sequence of elements
\[
i_{\beta} \;\in\; J_{\beta} \qquad\quad (\, \beta \;\leq\; \alpha \,) \; .
\]
If $i_{\beta}$ is constructed, since $J_{\beta+1}$ is cofinal, there exists $i_{\beta+1} \geq i_{\beta}$ in $J_{\beta+1}$, O.K. If $\beta$ is a limit ordinal, we take
\[
i_{\beta} \;=\; \sup_{\beta'<\beta} i_{\beta'} \; ,
\]
$i_{\beta}$ is in each $I_{\beta_{0}}$ for $\beta_{0} < \beta$,
\msp{page 72}
since the $i_{\beta'}$ for $\beta_{0} \leq \beta' < \beta$ are all in $J_{\beta_{0}} \supset J_{\beta'}$, hence also their limit, hence
\[
i_{\beta} \;\in\; \bigcap_{\beta_{0}<\beta} J_{\beta_{0}} \;=\; J_{\beta} \; .
\]
So transfinite induction works, and thus gives us an element $s_{\alpha}$ of $J_{\alpha}$\,!

I am also going to prove in due form the assertion for $\bigcup_{\alpha\in A} M_{\alpha}$, when $A$ is finite. Setting this time $M = \bigcup_{\alpha} M_{\alpha}$, we must prove, using the relations $(*)$ (p. 68), that we have
\[
\begin{aligned}
&X \in M &&\Longleftrightarrow \quad &&\text{the set } \ I_{M} \ \ (= \ \textstyle\bigcup_{\alpha} I_{M_{\alpha}})\\
&(\textit{i.e.} \ \exists \ \alpha \text{ with } X \in M_{\alpha}) && &&\text{is cofinal in } I \ ,
\end{aligned}
\tag*{$(**)$}
\]
and using $(*)$ the first term of this equivalence to be established means that there exists $\alpha$ with $I_{M_{\alpha}}$ cofinal in $I$. Of course, if one of the $I_{M_{\alpha}}$ is cofinal in $I$, a fortiori their union. But conversely, if the union is cofinal, it is filtering and a union of the (filtering) subsets $I_{M_{\alpha}}$, one of which must therefore be cofinal in it by Lemma 5.14, hence also in $I$, \textit{q.e.d.}

\msp{page 73}
\textbf{Corollary 5.16} (of Proposition 5.13). \textit{Every finite union of L-accessible subsets of $\mathcal{M}$ is L-accessible. Every intersection $\bigcap M_{\alpha}$ of a small family of L-accessible subsets of $\mathcal{M}$, is L-accessible.}

\printedpage{41}
% head: A. Grothendieck : Dérivateurs, Ch. XV, Édition des Universités de Montpellier II et Paris VII
% foot: Version 5 novembre 2010

\section*{6\quad Abortive attempt at a generation criterion for L-accessible subsets of $\Fl(\mathcal{M})$.}

Let $\mathcal{M}$ be a category equipped with a cardinal filtration $(\Filt^{\pi})_{\pi\geq\pi_{0}}$, and let $S$ be an object of Cat, let us consider
\[
\mathcal{M}(S) \;=\; \uHom(S,\mathcal{M}) \; ,
\]
and equip $\mathcal{M}(S)$ with the filtration by the
\[
\mathcal{M}^{\pi}(S) \;=\; \uHom(S,\mathcal{M}^{\pi} = \underline{\Filt}^{\pi}(\mathcal{M})) \; .
\]
I say that this is a cardinal filtration of $\mathcal{M}(S)$. We must check conditions $(a)$, $(b)$, $(c)$ of SGA 4 I 9.12. Conditions $(a)$ ($\Filt^{\pi}$ essentially small) and $(b)$ (condition $\mathrm{L}_{\pi_{0}}$, and stability of $\Filt^{\pi}$ under $\varinjlim\limits_{I}$ , $I$ filtering large relative to $\pi_{0}$ and of cardinality $\leq \pi$) pass at once from $\mathcal{M}$ to $\mathcal{M}(S)$. Condition $(c)$ is more delicate (representation
\msp{page 74}
of an object as a filtering $\varinjlim$ large relative to $\pi$ of objects of $\Filt^{\pi}$). I am not going to redo the work, done in a much more general and more delicate case in SGA 4 I §9, namely in Corollary 9.24 of Theorem 9.22 (p. 177). The only restriction is that there the objects of $\mathcal{M}$ are assumed accessible, and denoting by $\pi_{1}$ the smallest cardinal such that $\pi_{1} \geq \pi_{0}$ and that the objects of $\Filt^{\pi_{0}}$ are $\pi_{1}$-accessible, we restrict ourselves for the cardinal filtration of $\uHom(S,\mathcal{M})$ to the cardinals \fbox{$\pi \geq c = 2^{c_{0}}$,} where
\[
c_{0} \;=\; \sup(\pi_{1}, \card \Fl(S)) \; .
\]
This is a reasonable and innocuous restriction.

For the moment, I am interested above all in the case where $S = \boldsymbol{\Delta}_{1}$, hence in
\[
\uHom(\boldsymbol{\Delta}_{1},\mathcal{M}) \;=\; \uFl(\mathcal{M}) \; ,
\]
in which case we have $c_{0} = \pi_{1}$, so it suffices to restrict ourselves to the cardinals $\pi \geq 2^{\pi_{1}}$ to have a good cardinal filtration on $\uFl(\mathcal{M})$.

It is thus the set of the two conditions $(b)$, $(c)$ on $\mathcal{M}$ (p. 44) that is stable under passage to the $\uHom(S,\mathcal{M}) = \mathcal{M}(S)$. The same holds for condition $(a)$ (stability under small direct limits), and also for condition\msp{page 75}
$(d)$ (p. 45). Likewise the exactness condition $(e)$ (p. 46) on filtering $\varinjlim$ (they must transform monomorphisms into monomorphisms).

To be on the safe side, I will make on $\mathcal{M}$ the assumptions $(a)$ to $(e)$ inclusive (although $(e)$ seems a little artificial to me, we should be able to eliminate it, with a second wind\,!). These conditions thus have the desirable property of stability under passage to the $\mathcal{M}(S)$, and in particular under passage to $\Fl(\mathcal{M})$. Increasing $\pi_{0}$ if need be, we can therefore assume that the $\Fl^{\pi}(\mathcal{M}) = \mathrm{Hom}(\boldsymbol{\Delta}_{1},\mathcal{M}^{\pi})$ $(\pi \geq \pi_{0})$ also form a cardinal filtration of $\uFl(\mathcal{M})$.

We thus have at our disposal the notions of $\pi$-accessible, accessible, L-$\pi$-accessible, L-accessible subsets of $\Fl(\mathcal{M}) \ = \ \Ob\uFl(\mathcal{M})$. These notions are stable under finite unions and intersections, and the absolute notions (without $\pi$) stable under ``arbitrary'' intersections (\textit{i.e.} for \textit{small} families of subsets). This will be something particularly\printedpage{42}%
% head: A. Grothendieck : Dérivateurs, Ch. XV, Édition des Universités de Montpellier II et Paris VII
% foot: Version 5 novembre 2010
\msp{page 76}
crucial, the case of an intersection $C \cap W$\,!

The main result, for which I practiced the scales of the preceding section, is the following:

\setcounter{footnote}{25}
\textbf{Theorem 6.1}\footnote{Botched proof, statement perhaps \textit{false}\,!!}. \textit{Let $\mathcal{M}$ be a category satisfying the stated conditions:}
\begin{enumerate}[label={}]
\item[a)] \textit{Stability under small} $\varinjlim$ \textit{(not necessarily filtering).}
\item[b)] $\mathcal{M}$ \textit{equipped with a cardinal filtration} $(\Filt^{\pi})_{\pi\geq\pi_{0}}$.
\item[c)] \textit{The objects of $\mathcal{M}$ are accessible.}
\item[d)] \textit{The} $\Filt^{\pi}$ \textit{are stable under} $\varinjlim$ \textit{indexed by categories $I$ such that} $\card \Fl(I) \leq \pi$\textit{, and if $Z \in \Filt^{\pi}$, every quotient $Z'$ of $Z$ (such that $Z \longrightarrow Z'$ does not factor through a subobject $Z'_{0}$ of $Z'$, different from $Z'$) is in} $\Filt^{\pi}$.
\item[e)] \textit{Filtering direct limits in $\mathcal{M}$ transform monomorphisms into monomorphisms.}
\end{enumerate}
\textit{Let $C$ be an L-accessible subset of} $\Fl(\mathcal{M}) = \Ob\uFl(\mathcal{M})$, i.e. (5.12) \textit{a subset stable under} $\varinjlim$ \textit{that are filtering, and accessible. Suppose moreover that $C$ is made up of monomorphisms}\footnote{In the end, we get rid of this assumption.}\textit{, and that $C$ contains the} isomorphisms. \textit{Then there exists a small subset $C_{0} \subset C$, such that $C \subset \widetilde{C_{0}}$. More precisely, if $\pi$ is a cardinal such that $C$ is $\pi$-accessible, and if we set $C^{\pi} = C \cap \Fl^{\pi}(\mathcal{M})$, then we have}
\msp{page 77}
\[
C \;\subset\; \widetilde{C^{\pi}} \;=\; \widetilde{C_{0}} \; ,
\tag*{(6.1.1)}
\]
\textit{where $C_{0}$ is any small subset of $C^{\pi}$ such that every object of $C^{\pi}$ is isomorphic to an object of $C_{0}$.}

\textbf{Corollary 6.2.} \textit{We have $C_{*} = C_{0\,*}$, $\widetilde{C} = \widetilde{C_{0}}$.}

\textbf{Corollary 6.3.} \textit{Equivalent conditions on the L-accessible subset $C$ of} $\Fl(\mathcal{M})_{\mathrm{mono}}$\textit{, containing the isomorphisms:}
\begin{enumerate}[label={}]
\item[a)] \textit{$C = \widetilde{C}$.}
\item[b)] \textit{$C$ is stable under cobase change and under composition.}
\item[b$'$)] \textit{$C$ is stable under cobase change, and under simple amalgamated sums,} i.e. \textit{if $A \longrightarrow B_{1}$, $A \longrightarrow B_{2}$ are in $C$, then $A \longrightarrow B_{1} \amalg_{A} B_{2}$ as well.}
\item[c)] \textit{$C^{\pi}$ is stable in} $\Fl^{\pi}(\mathcal{M}) = \Fl(\mathcal{M}^{\pi})$ \textit{under cobase change and under composition.}
\item[c$'$)] \textit{$C^{\pi}$ is stable in} $\Fl(\mathcal{M}^{\pi})$ \textit{under cobase change and under amalgamated sums} (cf. b$'$).
\end{enumerate}
\textit{When these conditions are satisfied, then $(C,C_{*})$ is a Quillen couple, which falls under Theorem} 1 \textit{of section} 1 (p. 10).

The last assertion follows from Theorem 6.1 and from the corollary to Theorem 1 of section 1. We have
\printedpage{43}%
% head: A. Grothendieck : Dérivateurs, Ch. XV, Édition des Universités de Montpellier II et Paris VII
% foot: Version 5 novembre 2010
\msp{page 78}
on the other hand tautologically
\[
\begin{tikzcd}
(a) \arrow[r, Rightarrow] & (b) \arrow[r, Rightarrow] \arrow[d, Rightarrow] & (b') \arrow[d, Rightarrow] \\
 & (c) \arrow[r, Rightarrow] & (c') \ ,
\end{tikzcd}
\]
and it remains to prove $(c') \Longrightarrow (a)$, which we will do below.

Let us then prove 6.1. In (6.1.1) the equality $\widetilde{C_{0}} = \widetilde{C^{\pi}}$ is tautological. So everything comes down to proving
\[
C \;\subset\; \widetilde{C^{\pi}} \; .
\]
One difficulty here is that we do \textit{not} know that $\widetilde{C^{\pi}}$ is stable under filtering $\varinjlim$ -- we know it only for $C$. Let then
\[
A \;\into\; B
\]
be an arrow in $C$, let us prove that it is in $\widetilde{C^{\pi}}$. For this, we are going to construct, by transfinite induction, a \textit{strictly} increasing transfinite sequence of subobjects $A_{\alpha}$ of $B$ containing $A = A_{0}$,
\[
\underbrace{A}_{=\,A_{0}} \;\; \overset{i_{\alpha}}{\into} \;\; A_{\alpha} \;\; \overset{i_{\beta,\alpha}}{\into} \;\; A_{\beta} \xrightarrow{\ j_{\beta}\ } B \; ,
\]
such that we have
\[
i_{\alpha} \;\in\; \widetilde{C^{\pi}} \; , \qquad\quad \text{more generally the } i_{\beta,\alpha} \;\in\; \widetilde{C^{\pi}} \; , \qquad\quad j_{\alpha} \;\in\; C \; ,
\]
and that we have
\[
A_{\alpha} \;=\; \varinjlim_{\beta<\alpha} A_{\beta} \qquad\quad \text{if } \alpha \text{ limit ordinal} \; .
\]

\msp{page 79}
This already specifies the recursive definition for $\alpha = \alpha_{0}$ a limit ordinal. Let us note that the $A_{\alpha_{0}}$ constructed in terms of the $A_{\beta}$ $(\beta < \alpha_{0})$ does satisfy the conditions that
\begin{enumerate}[label={}]
\item[\textit{a})] the $i_{\alpha_{0},\beta} : A_{\beta} \longrightarrow A_{\alpha_{0}}$ are in $\widetilde{C^{\pi}}$, for $\widetilde{C^{\pi}}$ is stable under ``strict ordinal direct limits'' (\textit{cf.} 1.19, p. 8),
\item[\textit{b})] $j_{\alpha_{0}} \, : \, A_{\alpha_{0}} \longrightarrow B$ is in $C$, for it is the filtering direct limit of the system of the $A_{\beta} \longrightarrow B$ $(\beta < \alpha_{0})$ which are in $C$, and $C$ is stable under these limits.
\end{enumerate}
It remains then to specify how we pass from $A_{\alpha}$ to $A_{\alpha+1}$. This comes down to proving the

\textbf{Lemma 6.4.} \textit{Let $f : A \into B$ be in $C$, such that $f$ is not an isomorphism. Then there exists a subobject $A_{1}$ of $B$,} distinct from $A$, \textit{such that $A \longrightarrow A_{1}$ is in $\widetilde{C^{\pi}}$, and $A_{1} \longrightarrow B$ in $C$.}

If this lemma is admitted, this means that the transfinite construction continues as long as we find $A_{\alpha}$ such that $A_{\alpha} \neq B$, \textit{i.e.} $A_{\alpha} \into B$ is not an isomorphism. But since the $A_{\alpha}$ are all different, if $\pi$ is the set of subobjects of $B$ (we know that it is a

\end{document}
